% Options for packages loaded elsewhere
% Options for packages loaded elsewhere
\PassOptionsToPackage{unicode}{hyperref}
\PassOptionsToPackage{hyphens}{url}
\PassOptionsToPackage{dvipsnames,svgnames,x11names}{xcolor}
%
\documentclass[
  11pt,
]{article}
\usepackage{xcolor}
\usepackage[margin=2.5cm]{geometry}
\usepackage{amsmath,amssymb}
\setcounter{secnumdepth}{5}
\usepackage{iftex}
\ifPDFTeX
  \usepackage[T1]{fontenc}
  \usepackage[utf8]{inputenc}
  \usepackage{textcomp} % provide euro and other symbols
\else % if luatex or xetex
  \usepackage{unicode-math} % this also loads fontspec
  \defaultfontfeatures{Scale=MatchLowercase}
  \defaultfontfeatures[\rmfamily]{Ligatures=TeX,Scale=1}
\fi
\usepackage{lmodern}
\ifPDFTeX\else
  % xetex/luatex font selection
\fi
% Use upquote if available, for straight quotes in verbatim environments
\IfFileExists{upquote.sty}{\usepackage{upquote}}{}
\IfFileExists{microtype.sty}{% use microtype if available
  \usepackage[]{microtype}
  \UseMicrotypeSet[protrusion]{basicmath} % disable protrusion for tt fonts
}{}
\usepackage{setspace}
\makeatletter
\@ifundefined{KOMAClassName}{% if non-KOMA class
  \IfFileExists{parskip.sty}{%
    \usepackage{parskip}
  }{% else
    \setlength{\parindent}{0pt}
    \setlength{\parskip}{6pt plus 2pt minus 1pt}}
}{% if KOMA class
  \KOMAoptions{parskip=half}}
\makeatother
% Make \paragraph and \subparagraph free-standing
\makeatletter
\ifx\paragraph\undefined\else
  \let\oldparagraph\paragraph
  \renewcommand{\paragraph}{
    \@ifstar
      \xxxParagraphStar
      \xxxParagraphNoStar
  }
  \newcommand{\xxxParagraphStar}[1]{\oldparagraph*{#1}\mbox{}}
  \newcommand{\xxxParagraphNoStar}[1]{\oldparagraph{#1}\mbox{}}
\fi
\ifx\subparagraph\undefined\else
  \let\oldsubparagraph\subparagraph
  \renewcommand{\subparagraph}{
    \@ifstar
      \xxxSubParagraphStar
      \xxxSubParagraphNoStar
  }
  \newcommand{\xxxSubParagraphStar}[1]{\oldsubparagraph*{#1}\mbox{}}
  \newcommand{\xxxSubParagraphNoStar}[1]{\oldsubparagraph{#1}\mbox{}}
\fi
\makeatother


\usepackage{longtable,booktabs,array}
\usepackage{calc} % for calculating minipage widths
% Correct order of tables after \paragraph or \subparagraph
\usepackage{etoolbox}
\makeatletter
\patchcmd\longtable{\par}{\if@noskipsec\mbox{}\fi\par}{}{}
\makeatother
% Allow footnotes in longtable head/foot
\IfFileExists{footnotehyper.sty}{\usepackage{footnotehyper}}{\usepackage{footnote}}
\makesavenoteenv{longtable}
\usepackage{graphicx}
\makeatletter
\newsavebox\pandoc@box
\newcommand*\pandocbounded[1]{% scales image to fit in text height/width
  \sbox\pandoc@box{#1}%
  \Gscale@div\@tempa{\textheight}{\dimexpr\ht\pandoc@box+\dp\pandoc@box\relax}%
  \Gscale@div\@tempb{\linewidth}{\wd\pandoc@box}%
  \ifdim\@tempb\p@<\@tempa\p@\let\@tempa\@tempb\fi% select the smaller of both
  \ifdim\@tempa\p@<\p@\scalebox{\@tempa}{\usebox\pandoc@box}%
  \else\usebox{\pandoc@box}%
  \fi%
}
% Set default figure placement to htbp
\def\fps@figure{htbp}
\makeatother


% definitions for citeproc citations
\NewDocumentCommand\citeproctext{}{}
\NewDocumentCommand\citeproc{mm}{%
  \begingroup\def\citeproctext{#2}\cite{#1}\endgroup}
\makeatletter
 % allow citations to break across lines
 \let\@cite@ofmt\@firstofone
 % avoid brackets around text for \cite:
 \def\@biblabel#1{}
 \def\@cite#1#2{{#1\if@tempswa , #2\fi}}
\makeatother
\newlength{\cslhangindent}
\setlength{\cslhangindent}{1.5em}
\newlength{\csllabelwidth}
\setlength{\csllabelwidth}{3em}
\newenvironment{CSLReferences}[2] % #1 hanging-indent, #2 entry-spacing
 {\begin{list}{}{%
  \setlength{\itemindent}{0pt}
  \setlength{\leftmargin}{0pt}
  \setlength{\parsep}{0pt}
  % turn on hanging indent if param 1 is 1
  \ifodd #1
   \setlength{\leftmargin}{\cslhangindent}
   \setlength{\itemindent}{-1\cslhangindent}
  \fi
  % set entry spacing
  \setlength{\itemsep}{#2\baselineskip}}}
 {\end{list}}
\usepackage{calc}
\newcommand{\CSLBlock}[1]{\hfill\break\parbox[t]{\linewidth}{\strut\ignorespaces#1\strut}}
\newcommand{\CSLLeftMargin}[1]{\parbox[t]{\csllabelwidth}{\strut#1\strut}}
\newcommand{\CSLRightInline}[1]{\parbox[t]{\linewidth - \csllabelwidth}{\strut#1\strut}}
\newcommand{\CSLIndent}[1]{\hspace{\cslhangindent}#1}



\setlength{\emergencystretch}{3em} % prevent overfull lines

\providecommand{\tightlist}{%
  \setlength{\itemsep}{0pt}\setlength{\parskip}{0pt}}



 


\usepackage{setspace}
\usepackage{titlesec}
\usepackage{titletoc}

\titleformat{\chapter}[display]
  {\normalfont\huge\bfseries\centering}
  {\chaptertitlename\ \thechapter}{20pt}{\Huge}
\titlespacing*{\chapter}{0pt}{0pt}{40pt}
  
  % Section formatting
\titleformat{\section}
  {\normalfont\Large\bfseries\raggedright}
  {\thesection}{1em}{}
\titlespacing*{\section}{0pt}{3.5ex plus 1ex minus .2ex}{2.3ex plus .2ex}
  
% Subsection formatting
\titleformat{\subsection}
  {\normalfont\large\bfseries\raggedright}
  {\thesubsection}{1em}{}
\titlespacing*{\subsection}{0pt}{3.25ex plus 1ex minus .2ex}{1.5ex plus .2ex}

% Subsubsection formatting
\titleformat{\subsubsection}
  {\normalfont\normalsize\bfseries\raggedright}
  {\thesubsubsection}{1em}{}
\titlespacing*{\subsubsection}{0pt}{3.25ex plus 1ex minus .2ex}{1.5ex plus .2ex}

% Paragraph formatting (if needed)
\titleformat{\paragraph}[runin]
  {\normalfont\normalsize\bfseries}
  {\theparagraph}{1em}{}
\titlespacing*{\paragraph}{0pt}{3.25ex plus 1ex minus .2ex}{1em}
\makeatletter
\@ifpackageloaded{caption}{}{\usepackage{caption}}
\AtBeginDocument{%
\ifdefined\contentsname
  \renewcommand*\contentsname{Table of contents}
\else
  \newcommand\contentsname{Table of contents}
\fi
\ifdefined\listfigurename
  \renewcommand*\listfigurename{List of Figures}
\else
  \newcommand\listfigurename{List of Figures}
\fi
\ifdefined\listtablename
  \renewcommand*\listtablename{List of Tables}
\else
  \newcommand\listtablename{List of Tables}
\fi
\ifdefined\figurename
  \renewcommand*\figurename{Figure}
\else
  \newcommand\figurename{Figure}
\fi
\ifdefined\tablename
  \renewcommand*\tablename{Table}
\else
  \newcommand\tablename{Table}
\fi
}
\@ifpackageloaded{float}{}{\usepackage{float}}
\floatstyle{ruled}
\@ifundefined{c@chapter}{\newfloat{codelisting}{h}{lop}}{\newfloat{codelisting}{h}{lop}[chapter]}
\floatname{codelisting}{Listing}
\newcommand*\listoflistings{\listof{codelisting}{List of Listings}}
\makeatother
\makeatletter
\makeatother
\makeatletter
\@ifpackageloaded{caption}{}{\usepackage{caption}}
\@ifpackageloaded{subcaption}{}{\usepackage{subcaption}}
\makeatother
\usepackage{bookmark}
\IfFileExists{xurl.sty}{\usepackage{xurl}}{} % add URL line breaks if available
\urlstyle{same}
\hypersetup{
  colorlinks=true,
  linkcolor={blue},
  filecolor={Maroon},
  citecolor={Blue},
  urlcolor={Blue},
  pdfcreator={LaTeX via pandoc}}


\author{}
\date{}
\begin{document}


\setstretch{1.2}
\section{Methods}\label{methods}

\subsection{Study design and
population}\label{study-design-and-population}

The OsteoLaus study includs a subgroup of the CoLaus \textbar{}
PsyCoLaus study participants, including 1474 postmenopausal women aged
50 to 80 years at baseline {[}\citeproc{ref-shevroja2019}{1}{]}. The
detailed description of the recruitment process and methodology for
CoLaus and OsteoLaus has been described elsewhere
{[}\citeproc{ref-firmann2008}{2}{]} {[}\citeproc{ref-shevroja2019}{1}{]}
and can be found at
\href{http://www.colaus-psycolaus.ch/}{www.colaus-psycolaus.ch}. In
summary, \ldots{}

CoLaus conducted at Centre Hospitalier Universitaire Vaudois

OsteoLaus (conducted at Centre for Bone Diseases at the LUMC) was
initiated to focus on bone health and is a population-based prospective
cohort from Lausanne, Switzerland. Participants were recruited between
September 2009 and September 2012 if they participated in the baseline
measurement of the CoLaus study. Women with the following inclusion
criteria were selected: written informed consent, BMI between 15
kg/m\textsuperscript{2} and 40 kg/m\textsuperscript{2}, at least two
vertebrae and the femur accessible for measurement, available
information on previous osteoporotic fractures, QUS of the heel
accessible for measurement, no known disease affecting bone metabolism,
and no known disorders affecting skeletal sites. On average,
measurements were taken every 2.5 years, including the baseline and
follow-ups. The timeline was as follows: F1: Follow-up 2009--2013, F2:
Follow-up 2014--2018, F3: Follow-up 2018--2021, F4: Follow-up
2022--2026.

\subsubsection{Exclusion Criteria}\label{exclusion-criteria}

Figure 1 illustrates the study population.

\subsection{Exposure}\label{exposure}

FFQ self-reported food questionnaire

\subsection{Covariates}\label{covariates}

Figure 2 shows the DAG.

\subsubsection{Sociodemographic}\label{sociodemographic}

Educational level was self-reported and expressed as the highest level
achieved. The levels are divided in three categories corresponding to
International Standard Classification of Education (ISCED) 0-2, ISCED 3,
and ISCED 4-8, which are labelled low, medium, and high, respectively
{[}\citeproc{ref-ISCED}{3}{]}.

\subsubsection{Lifestyle}\label{lifestyle}

Smoking has been recognised as a risk factor for muscle wasting
\emph{{[}SOURCE:}
\href{https://doi.org/10.1164/rccm.201410-1830PP}{10.1164/rccm.201410-1830PP}\emph{{]}.}
Participants' smoking history was assessed via self-report. Individuals
were categorized as current smokers if they reported active tobacco use
at the time of the study. Those who reported a history of smoking but
had ceased use prior to the study were classified as former smokers.
Participants with no history of tobacco initiation were defined as never
smokers.

Alcohol intake was assessed using a self-reported questionnaire covering
consumption patterns over the previous 7 days. Alcohol consumption was
analyzed as both a categorical and a continuous variable. Following
Swiss national guidelines, a standard drink was defined as a single
glass of wine, a bottle of beer, or a shot of spirits, each containing
approximately 10 g to 12 g of pure ethanol
{[}\citeproc{ref-standardglas}{4}{]}. Participants were categorized into
four groups based on their total weekly intake to reflect sex-specific
physiological thresholds:

\begin{itemize}
\tightlist
\item
  Non-drinkers: 0 g/day (0 units/week).
\item
  Light drinkers: \textgreater0 to 6 g/day (1--3 units/week),
  representing occasional consumption.
\item
  Moderate drinkers: \textgreater6 to 12 g/day (4--7 units/week),
  representing frequent consumption within the Swiss low-risk limit for
  women.
\item
  Heavy drinkers: \textgreater12 g/day (\textgreater7 units/week),
  categorized as hazardous consumption exceeding the Swiss national
  guidelines for women.
\end{itemize}

While this categorization provides a structured metric for analysis, it
is noted that self-reported dietary and alcohol questionnaires often
slightly underestimate absolute intake levels
{[}\citeproc{ref-kerr2012}{5},\citeproc{ref-morabia1994}{6}{]}.
Validation studies in the Swiss population (Geneva) have specifically
demonstrated that semi-quantitative food frequency questionnaires (FFQs)
tend to underestimate absolute alcohol consumption when compared to more
intensive 24-hour recall methods {[}\citeproc{ref-morabia1994}{6}{]}.

Physical activity

\subsubsection{Dietary assessment}\label{dietary-assessment}

Total energy intake

Vitamin D

\subsubsection{Clinical}\label{clinical}

The relationship between \textbf{diabetes mellitus} and muscle mass is
complex and varies substantially by diabetes type and patient
characteristics. Type 1 diabetes demonstrates the most pronounced
effects, with untreated patients experiencing dramatic muscle loss and
experimental models showing 20-30\% decreases in muscle mass and
myofiber area \emph{D. Sala et al., 2015}. In contrast, the effects of
type 2 diabetes on muscle mass are more heterogeneous, with most
patients showing no changes or even increases in muscle mass \emph{D.
Sala et al., 2015}. However, longitudinal studies reveal that certain
populations, particularly older adults, experience accelerated muscle
mass decline compared to non-diabetic controls \emph{A. Koster et al.,
2015.}

The relationship between type 2 diabetes and sarcopenia is bidirectional
and multifactorial {[}\citeproc{ref-chen2023}{7}{]}. While diabetes can
promote muscle deterioration through insulin resistance, chronic
inflammation, oxidative stress, and hyperglycemia-induced protein
breakdown pathways \emph{Y. Hirata et al., 2019}, reduced muscle mass
conversely increases diabetes risk and worsens glucose metabolism
\emph{D. Sizoo et al., 2022.} Clinical evidence demonstrates that muscle
deterioration in type 2 diabetes patients occurs independently of
disease duration, metabolic control, and microvascular complications,
suggesting intrinsic pathophysiological mechanisms beyond traditional
diabetic complications \emph{N. Guerrero et al., 2016.} Importantly,
therapeutic interventions including glycemic control optimization and
insulin therapy can improve muscle mass outcomes, indicating the
potential reversibility of diabetes-associated muscle changes \emph{K.
Sugimoto et al., 2020.}

Based on these findings, patients are allocated in the following 4
categories:

\begin{itemize}
\item
  \textbf{No diabetes} (negative on all diagnostic criteria)
\item
  \textbf{Diabetes, no insulin} (positive on FPG/HbA1c criteria, oral
  medications only)
\item
  \textbf{Diabetes with insulin} (insulin use regardless of other
  treatments)
\item
  \textbf{IGT/prediabetes} (IGT positive but not meeting full diabetes
  criteria)
\end{itemize}

\subsubsection{\texorpdfstring{\textbf{Anthropometric}}{Anthropometric}}\label{anthropometric}

All participants had their height measured using the same portable
stadiometer (Seca version 216, Seca, Chino, CA, USA) with precision 0.1
cm, and body weight with the same electronic scale (Seca Clara 803,
Seca, Chino, CA, USA) with precision 0.1 kg. Both measurements were
performed with the participant barefoot and in minimum clothing. Body
mass index (BMI) was calculated as weight/height\textsuperscript{2}. BMI
was categorized into underweight (\textless{} 18.5
kg/m\textsuperscript{2}), normal (18.5-24.9 kg/m\textsuperscript{2}),
overweight (25.0--29.9 kg/m\textsuperscript{2}), and obese (≥ 30.0
kg/m\textsuperscript{2}).

\subsection{Outcome Variables}\label{outcome-variables}

Handgrip strength (kg), measured using a JAMAR dynamometer

Appendicular lean mass (ALM) and derived indices (ALMI and ALM/BMI for
sensitivity analysis), measured by DXA (GE Lunar iDXA™)

Gait speed (m/s), assessed using a 6-meter walk test

\subsection{Statistical Methods}\label{statistical-methods}

Participants were classified into quartiles according to their baseline
daily intake of dairy products (grams per day) for descriptive analysis.

\subsection{Ethics}\label{ethics}

xx

\section*{References}\label{bibliography}
\addcontentsline{toc}{section}{References}

\phantomsection\label{refs}
\begin{CSLReferences}{0}{1}
\bibitem[\citeproctext]{ref-shevroja2019}
\CSLLeftMargin{1. }%
\CSLRightInline{Enisa Shevroja, Pedro Marques-Vidal, Bérengère
Aubry-Rozier, Gabriel Hans, Fernando Rivadeneira, Olivier Lamy, Didier
Hans. Cohort {Profile}: {The OsteoLaus} study. International Journal of
Epidemiology. 2019 Aug;48(4):1046--1047g.
doi:\href{https://doi.org/10.1093/ije/dyy276}{10.1093/ije/dyy276}}

\bibitem[\citeproctext]{ref-firmann2008}
\CSLLeftMargin{2. }%
\CSLRightInline{Mathieu Firmann, Vladimir Mayor, Pedro Marques Vidal,
Murielle Bochud, Alain Pécoud, Daniel Hayoz, Fred Paccaud, Martin
Preisig, Kijoung S. Song, Xin Yuan, Theodore M. Danoff, Heide A.
Stirnadel, Dawn Waterworth, Vincent Mooser, Gérard Waeber, Peter
Vollenweider. The {CoLaus} study: A population-based study to
investigate the epidemiology and genetic determinants of cardiovascular
risk factors and metabolic syndrome. BMC Cardiovascular Disorders. 2008
Mar;8(1):6.
doi:\href{https://doi.org/10.1186/1471-2261-8-6}{10.1186/1471-2261-8-6}}

\bibitem[\citeproctext]{ref-ISCED}
\CSLLeftMargin{3. }%
\CSLRightInline{International {Standard Classification} of {Education}
({ISCED}).
https://ec.europa.eu/eurostat/statistics-explained/index.php?title=International\_Standard\_Classification\_of\_Education\_(ISCED).}

\bibitem[\citeproctext]{ref-standardglas}
\CSLLeftMargin{4. }%
\CSLRightInline{Das {Standardglas Alkohol}... - {Alcohol Facts}.
https://www.aktionstag-alkoholprobleme.ch/de/das-standardglas-alkohol.}

\bibitem[\citeproctext]{ref-kerr2012}
\CSLLeftMargin{5. }%
\CSLRightInline{William C. Kerr, Tim Stockwell. Understanding standard
drinks and drinking guidelines. Drug and Alcohol Review.
2012;31(2):200--5.
doi:\href{https://doi.org/10.1111/j.1465-3362.2011.00374.x}{10.1111/j.1465-3362.2011.00374.x}}

\bibitem[\citeproctext]{ref-morabia1994}
\CSLLeftMargin{6. }%
\CSLRightInline{Alfredo Morabia, Martine Bernstein, Shiriki Kumanyika,
Ann Sorenson, Isabelle Mabiala, Bettina Prodolliet, Irène Rolfo, Ba Lau
Luong. {Développement et validation d'un questionnaire alimentaire
semi-quantitatif à partir d'une enquête de population}. Sozial- und
Präventivmedizin. 1994 Nov;39(6):345--69.
doi:\href{https://doi.org/10.1007/BF01299666}{10.1007/BF01299666}}

\bibitem[\citeproctext]{ref-chen2023}
\CSLLeftMargin{7. }%
\CSLRightInline{Huiling Chen, Xiaojing Huang, Meiyuan Dong, Song Wen,
Ligang Zhou, Xinlu Yuan. The {Association Between Sarcopenia} and
{Diabetes}: {From Pathophysiology Mechanism} to {Therapeutic Strategy}.
Diabetes, Metabolic Syndrome and Obesity. 2023 May;16:1541--54.
doi:\href{https://doi.org/10.2147/DMSO.S410834}{10.2147/DMSO.S410834}}

\end{CSLReferences}




\end{document}
